\documentclass[11pt,a4paper]{article}
\usepackage[T1]{fontenc}
\usepackage{isabelle,isabellesym}
\usepackage{amsmath,amssymb}

% This should be the last package used.
\usepackage{pdfsetup}

\urlstyle{rm}
\isabellestyle{it}

\begin{document}

\title{Stokes' Theorem on Smooth Manifolds with Boundary}
\author{Lawrence C.\ Paulson}
\maketitle

\begin{abstract}
We formalise Stokes' theorem for oriented smooth manifolds with boundary
in Isabelle/HOL. The main result equates the integral of the exterior
derivative of a compactly supported near-top form with the integral of
its boundary restriction. Forms are represented by compatible
coordinate families, and both integrals are defined using finite
partitions of unity. The development establishes the local half-space
identity, the induced boundary orientation, chart independence of the
integrals, and the global gluing theorem. It also treats dimension one
separately and gives a two-chart instance for the closed unit disk.
\end{abstract}

\section{Scope}

For positive dimension, the boundary is modelled on a closed
half-space. In dimensions at least two, the integrated forms have
degree one less than the manifold dimension, represented by their
coordinate coefficient fields. The one-dimensional development uses
an exact-form interface and a chart-independent one-form integral.
The disk instance identifies its boundary with the unit circle.
The Green-form calculation checks the usual
$\partial_x Q-\partial_y P$ sign in local coordinates; an ordinary
area-and-contour-integral Green theorem is not claimed here.

\section{Formal development}

\input{session}

\end{document}
